% GENERATED FILE. Edit canonical lessons, then run npm run generate:books.
% canonical-source-hash: fnv1a64:d138302d82dd7cf3
% canonical-lessons: PT-C260-no-fim, PT-C260-contar, PT-C260-leitura, PT-C260-revisao-historia

\chapter{First, Then, In the End --- A Story Told in Order}
\label{ch:pt-a-story-told-in-order}
\hlchaptermodality{260}

\begin{chapteropening}
\textbf{By the end of this chapter:} \emph{I can tell a short connected story in order with primeiro, depois and no fim, and read one of about a hundred words.}

Complete PT-C260-revisao-historia: retell the reading's journey aloud, choosing scene or event for every verb.
\end{chapteropening}

\section[primeiro, depois, no fim]{primeiro, depois, no fim — the events, in order}

\label{lesson:PT-C260-no-fim}



\begin{quote}
Say \textquotedblleft{}first\textquotedblright{}, the ordinal from chapter twenty-seven. (\emph{Primeiro.}) Then \textquotedblleft{}afterwards\textquotedblright{}. (\emph{Depois.})

Now the same words go to work on a story. The scene is painted; the events need a running order.
\end{quote}

\begin{grammarlens}[title={Grammar Lens: three places in a story}]
\begin{quote}
\textbf{Primeiro} comprei o bilhete. — First I bought the ticket. \textbf{Depois} esperei o comboio. — Then I waited for the train. \textbf{No fim}, voltei para casa. — In the end, I went back home.
\end{quote}

\textbf{Depois} is the plain \textquotedblleft{}then, after that\textquotedblright{}. It stands between two events and points from the first to the second, and you can use it as many times as the story has events. \textbf{A seguir}, which you met as \textquotedblleft{}next\textquotedblright{}, does the same job.

\textbf{Depois} can also take a noun, with \textbf{de}: \emph{depois do jantar}, \textquotedblleft{}after dinner\textquotedblright{}, where \emph{de} and \emph{o} weld into \textbf{do}. \emph{A seguir} takes \emph{a} instead, \emph{a seguir ao jantar}, and often stands alone.

\textbf{No fim} is \textquotedblleft{}in the end\textquotedblright{}: the last event, and usually the outcome. You have \textbf{finalmente} too, \textquotedblleft{}finally, at last\textquotedblright{}, and it carries relief: something you were waiting for has come. \emph{No fim} simply closes the story. \emph{Fazia frio e chovia, e no fim voltei para casa contente.}

Notice what these words do \textbf{not} touch: the scene. \emph{Fazia frio} needs no \emph{depois}. Order belongs to the events.
\end{grammarlens}

\begin{cousinweb}
\textbf{depois} is usually traced to Latin \emph{dē} and a spoken Latin \emph{postius}, a form of \emph{post}, \textquotedblleft{}after, behind\textquotedblright{} (Lewis \& Short, under \emph{post}); EtymDB 2.1 follows it only to Old Portuguese \emph{depois}. English keeps \emph{post} in \textbf{postpone} and \textbf{postscript}.

\textbf{o fim}, \textquotedblleft{}the end\textquotedblright{}, has no parent in EtymDB 2.1. It is usually traced to Latin \emph{fīnis}, \textquotedblleft{}boundary, end\textquotedblright{} (Lewis \& Short), which English kept in \textbf{final} and \textbf{finish}, and which also gave \textbf{finalmente}.
\end{cousinweb}

\subsection*{Your turn}
Say these aloud:

\begin{itemize}
\raggedright
  \item \emph{Primeiro comprei o bilhete. Depois esperei o comboio.}
  \item \emph{depois do jantar} — then \emph{a seguir}, on its own
  \item \emph{No fim, voltei para casa.}
  \item a scene, then events — \emph{Fazia frio. Primeiro esperei. Depois, quando o comboio chegou, entrei.}
\end{itemize}

\begin{tcolorbox}[breakable,colback=teal!4,colframe=teal!35!black,title={Before you move on}]
Which of \emph{depois} and \emph{a seguir} takes a noun with \emph{de}? (\emph{\textbf{Depois}} — \emph{depois do jantar}.) Which closes a story, and which says \textquotedblleft{}at last\textquotedblright{}? (\emph{\textbf{No fim}} closes it; \emph{\textbf{finalmente}} is \textquotedblleft{}at last\textquotedblright{}.) Which Latin word is inside \emph{depois}? (\emph{\textbf{Post}}, \textquotedblleft{}after\textquotedblright{}.)
\end{tcolorbox}
\section[contar uma história]{contar uma história — to tell a story, which is to count it}

\label{lesson:PT-C260-contar}



\begin{quote}
You can hold a scene still, put events in order, and say what a reason led to with \emph{por isso}. This lesson adds one small word, then puts it all together.
\end{quote}

\begin{grammarlens}[title={Grammar Lens: mas, the turn in a story}]
\textbf{mas} means \textquotedblleft{}but\textquotedblright{}. It joins two facts that pull against each other:

\begin{quote}
\emph{Estava sol, \textbf{mas} fazia frio.} — The sun was out, but it was cold. \emph{Esperei muito, \textbf{mas} o comboio não chegou.} — I waited a long time, but the train did not come.
\end{quote}

In a story it marks the turn. \emph{E} adds; \emph{mas} turns.
\end{grammarlens}

\begin{grammarlens}[title={Grammar Lens: the backbone of a told story}]
\textbf{contar uma história} is \textquotedblleft{}to tell a story\textquotedblright{}, with a \emph{contar} you already have. Here is what it tells, with both pasts to work with:

\begin{quote}
\textbf{Era} sábado. \textbf{Fazia} calor e \textbf{estava} sol. \textbf{Primeiro fui} ao mercado. \textbf{Depois encontrei} um amigo, e \textbf{falámos} muito. \textbf{Mas no fim começou} a chover, \textbf{por isso voltei} para casa.
\end{quote}

\begin{quote}
\emph{It was Saturday. It was hot and the sun was out. First I went to the market. Then I met a friend, and we talked a lot. But in the end it started to rain, so I went back home.}
\end{quote}

\emph{Era}, \emph{fazia}, \emph{estava}: the scene, held still. \emph{Fui}, \emph{encontrei}, \emph{falámos}, \emph{começou}, \emph{voltei}: the events, put in order by \emph{primeiro}, \emph{depois} and \emph{no fim}. \emph{Mas} marks the turn; \emph{por isso} says what the rain led to.
\end{grammarlens}

\begin{cousinweb}
\textbf{contar} is inherited from Latin \emph{computāre}, \textquotedblleft{}to reckon up\textquotedblright{} (EtymDB 2.1), which is \emph{com-}, \textquotedblleft{}together\textquotedblright{}, plus \emph{putāre}, \textquotedblleft{}to settle an account\textquotedblright{} (Lewis \& Short). Portuguese kept one verb for both jobs: \emph{contar até dez}, to count to ten, and \emph{contar uma história}. English \textbf{recount} has the same picture.

\textbf{mas} is inherited from Latin \emph{magis}, \textquotedblleft{}more\textquotedblright{} (EtymDB 2.1), the same Latin word that gave \textbf{mais}. \textquotedblleft{}More\textquotedblright{} became \textquotedblleft{}rather\textquotedblright{}, and \textquotedblleft{}rather\textquotedblright{} became \textquotedblleft{}but\textquotedblright{}.
\end{cousinweb}

\subsection*{Your turn}
Say these aloud:

\begin{itemize}
\raggedright
  \item the scene — \emph{Era sábado. Fazia calor e estava sol.}
  \item the events — \emph{Primeiro fui ao mercado. Depois encontrei um amigo.}
  \item the turn and the end — \emph{Mas no fim começou a chover, por isso voltei para casa.}
  \item \emph{contar até dez}, then \emph{contar uma história}
\end{itemize}

\begin{tcolorbox}[breakable,colback=teal!4,colframe=teal!35!black,title={Before you move on}]
In your story, which past held the sunshine and which one took you to the market? (\emph{\textbf{Estava}} — imperfeito; \emph{\textbf{fui}} — perfeito.) Which word marked the turn? (\emph{\textbf{Mas.}}) What are the two jobs of \emph{contar}? (\textbf{Counting} and \textbf{telling}, both from Latin \emph{computāre}.)
\end{tcolorbox}
\section[(reading)]{(reading) — a journey, told in both pasts}

\label{lesson:PT-C260-leitura}



\begin{quote}
Chapter twenty-nine gave you forty-two words in a row, all in the present. This passage is twice as long, and it happened yesterday.

Every word in it was taught before this page.
\end{quote}

\begin{tcolorbox}[breakable,skin=enhanced,colback=orange!6,colframe=orange!45!black,boxrule=0.5pt,arc=1mm,left=6pt,right=6pt,top=4pt,bottom=4pt,fonttitle=\bfseries,title={Read this}]
\begin{quote}
Ontem era sábado. De manhã fazia frio e chovia. Eu estava na estação com o meu amigo. Havia muitas pessoas, mas o comboio não chegava. Primeiro esperámos no cais. Tínhamos fome, por isso depois tomámos um café e comemos pão. Quando o comboio chegou, entrámos.
\end{quote}

\begin{quote}
A viagem era longa, mas não era aborrecida: o meu amigo contava histórias, e eu ouvia. Lá fora ainda chovia, mas no comboio estava quente. No fim, à noite, chegámos a Lisboa. O céu estava limpo e havia estrelas. Eu estava contente, e o meu amigo também.
\end{quote}

Read it once without stopping. Do not translate as you go.

Then read it a second time and sort every verb into one of two piles: the scene, or what happened.
\end{tcolorbox}

\begin{culture}
The events pile is short: \emph{esperámos}, \emph{tomámos}, \emph{comemos}, \emph{chegou}, \emph{entrámos}, \emph{chegámos}. Six things happened, and \emph{primeiro}, \emph{depois} and \emph{no fim} put them in order. \emph{Comemos} looks like the present, and it is spelled the same; the story around it tells you it is past.

Everything else is scene: the cold, the rain, the crowd, the hunger, the stars. A story is mostly background.

\emph{O comboio não chegava} is scene: the not-arriving went on. \emph{Quando o comboio chegou} is the event that ended it. And \emph{contava} is scene too; \emph{contou} would have made it one finished event.
\end{culture}

\subsection*{Your turn}
\begin{itemize}
\raggedright
  \item \emph{Say it:} the passage aloud, once, without stopping
\end{itemize}

\begin{tcolorbox}[breakable,colback=teal!4,colframe=teal!35!black,title={Before you move on}]
How many words in that passage were new? (\textbf{None.})

Why is \emph{contava} in the imperfeito? (The telling \textbf{went on} while the train moved: it is scene.)

Which word opened the line where a scene ended and an event began? (\emph{\textbf{Quando.}})
\end{tcolorbox}
\section[(review)]{Review — tell it back}

\label{lesson:PT-C260-revisao-historia}



\begin{quote}
Say \textquotedblleft{}first\textquotedblright{}, \textquotedblleft{}then\textquotedblright{} and \textquotedblleft{}in the end\textquotedblright{}. (\emph{Primeiro, depois, no fim.})
\end{quote}

\begin{grammarlens}[title={What you've built}]
You have read a story. Now you tell one, which is the same skill turned round: a reader sorts verbs into scene and event, and a teller has to choose them.

The journey from the reading, in five cues. Answer each one in Portuguese, out loud, in a full sentence. Choose the past as you go: is this cue the scene, or something that happened? And keep one \emph{mas} for the turn.
\end{grammarlens}

\subsection*{Your turn}
Say these aloud:

\begin{itemize}
\raggedright
  \item the weather that morning (\emph{De manhã fazia frio e chovia.})
  \item where you were, and what the train was not doing (\emph{Eu estava na estação, mas o comboio não chegava.})
  \item the first thing you did, then the second (\emph{Primeiro esperámos no cais. Depois tomámos um café.})
  \item what your friend was doing on the train (\emph{O meu amigo estava a contar histórias.})
  \item where you were in the end, and how you felt (\emph{No fim chegámos a Lisboa. Eu estava contente.})
\end{itemize}

Now tell the whole journey once more, in your own order and your own words, without the cues.

\begin{tcolorbox}[breakable,colback=teal!4,colframe=teal!35!black,title={Before you move on}]
Which of your five cues were scene? (The weather, where you were, the friend's telling, how you felt.) Which were events? (Waiting, the coffee, arriving.) Which Latin word is behind \emph{depois}? (\emph{\textbf{Post}}.) And what are you doing, word for word, when you \emph{contar uma história}? (\textbf{Counting it} — from \emph{computāre}.)
\end{tcolorbox}
